% NEW ACRONYM — add to formatting/Acronyms.tex:
% \DeclareAcronym{POMDP}{short=POMDP, long=Partially Observable Markov Decision Process, tag=abbrev, sort=POMDP}

% Content moved from Background §2.5 "Policy Learning Pipeline" to Methodology,
% because it describes the specific MDP instantiation and IsaacLab workflow
% used in this project's tasks. General RL theory remains in Background §2.1.
\label{sec:policy_learning}

In this thesis, \ac{PPO} is used as the primary learning algorithm, together with \ac{GAE} for advantage estimation, following standard actor--critic practice~\cite{Schulman2015GAE,Schulman2017PPO}. The training loop is implemented with the \texttt{skrl} library, which provides modular \ac{PPO} components (policy, value function, rollout storage, and optimization routines) and supports integration with simulator-based environments~\cite{SerranoMunoz2023skrl}. This allows the thesis to focus on environment formulation (observation design, action parameterization, reward shaping, termination conditions, and robustness measures such as randomization) while relying on a consistent and reproducible \ac{PPO} implementation~\cite{SerranoMunoz2023skrl}.

Reinforcement learning-based control is typically realized as an iterative interaction loop between an agent (policy) and an environment (simulation), in which behavior is improved by maximizing an expected return defined over episodes~\cite{SuttonBarto2018,Puterman1994}. In IsaacLab, task construction is formalized through an explicit \ac{MDP} specification and a software architecture that decomposes environment logic into reusable terms for observations, actions, rewards, terminations, events, and curriculum scheduling~\cite{nvidia_isaaclab_mdp_managers_getting_started,isaaclab_task_design_workflows,isaaclab_mdp_terms_api}. This section describes the \ac{RL} pipeline under this workflow and derives an \ac{MDP} formulation suitable for the reach and place tasks in simulation.